\section*{Lab 1: Measuring the Acceleration of a Cart on a Track}

\subsection*{1. Objective}
To measure the linear acceleration of a cart moving along a low-friction track and verify the relationship between force, mass, and acceleration.

\subsection*{2. Apparatus}
\begin{itemize}
    \item Low-friction track
    \item Dynamics cart
    \item Motion sensor or photogates with timer
    \item Pulley and string
    \item Hanging mass (20--50 g)
    \item Meter scale
    \item Computer or data-logger (e.g., Logger Pro, Pasco Capstone)
\end{itemize}

\subsection*{3. Theory}
According to Newton's Second Law,
\[
F = ma.
\]
If a hanging mass provides the force on the cart, then the net force is
\[
F = m_h g,
\]
where \(m_h\) is the hanging mass and \(g\) is the acceleration due to gravity.

The theoretical acceleration of the system is
\[
a_{\text{theory}} = \frac{m_h g}{m_h + m_c},
\]
where \(m_c\) is the mass of the cart.

\subsection*{4. Procedure}
\begin{enumerate}
    \item Set up the track horizontally and ensure minimal friction.
    \item Attach a string to the cart, pass it over a pulley, and connect the hanging mass.
    \item Position the motion sensor so it can record the cart's motion.
    \item Measure and record the masses \(m_c\) and \(m_h\).
    \item Raise the cart slightly so the hanging mass is lifted; release the cart without applying a push.
    \item Record the motion using the sensor or timer system.
    \item Generate a velocity--time graph from the recorded data.
    \item Determine the experimental acceleration \(a_{\text{exp}}\) from the slope of the velocity--time graph.
    \item Repeat the measurement at least three times for consistency.
\end{enumerate}

\subsection*{5. Data and Calculations}
From the velocity--time graph, the slope gives the experimental acceleration:
\[
a_{\text{exp}} = \frac{\Delta v}{\Delta t}.
\]

The theoretical acceleration is calculated using
\[
a_{\text{theory}} = \frac{m_h g}{m_h + m_c}.
\]

Percentage error:
\[
\% \text{Error} = 
\left(
\frac{|a_{\text{exp}} - a_{\text{theory}}|}{a_{\text{theory}}}
\right) \times 100.
\]

\subsection*{6. Conclusion}
The conclusion should state whether the experimental value of acceleration agrees with the theoretical value and discuss possible sources of error, including:
\begin{itemize}
    \item Friction in the track or pulley
    \item Timing or sensor inaccuracies
    \item Misalignment of the motion detector
    \item Variations in release technique
\end{itemize}
